\documentclass[12pt]{article}
\usepackage[utf8]{inputenc}
\usepackage[T1]{fontenc}
\usepackage[margin=1in]{geometry}
\usepackage{setspace}
\usepackage{xcolor}
\usepackage[colorlinks=true, urlcolor=blue!60!black, linkcolor=blue!60!black]{hyperref}
\setlength{\parindent}{0em}
\setlength{\parskip}{0.6em}
\renewcommand{\baselinestretch}{1.0}
\begin{document}
\begin{center}
{\LARGE \textit{Michael H. v. Gerald D.}, 491 U.S. 110 (1989)\par}
\vspace{0.5em}
{\large Case Brief --- Two-Page Edition\par}
\end{center}

\section{Memory Jogger}
Upholding California's conclusive presumption of marital paternity against a biological father's due process claim, the plurality's Footnote 6 launched the debate over defining rights at ``the most specific level'' of tradition (\href{https://supreme.justia.com/cases/federal/us/491/110/}{491 U.S. 110}).

\section{Detailed Case Facts}
Victoria D. was born in 1981 to Carole D., who was married to Gerald D.; blood tests showed a 98.07\% probability that Carole's lover, Michael H., was the father. Michael lived with Victoria at times and supported her, but Carole reconciled with Gerald in 1984. California Evidence Code \S~621 conclusively presumed a cohabiting wife's child to be the husband's.

\section{Procedural History}
Michael's filiation action ended in summary judgment for Gerald under \S~621; the Court of Appeal affirmed and the California Supreme Court denied review.

\section{Judicial Votes}
Affirmed, 5--4, without a majority opinion: Scalia wrote for Rehnquist, O'Connor, and Kennedy (only Rehnquist joined Footnote 6); Stevens concurred in the judgment; Brennan (with Marshall and Blackmun) and White (with Brennan) dissented.

\section{Holding}
A State may apply a conclusive marital presumption to bar a biological father's paternity claim to a child born into an extant, cohabiting marriage whose spouses wish to raise the child. Footnote 6 drew only two votes.

\section{Analysis of Opinions}
Scalia treated \S~621 as a substantive rule and required that a liberty interest be ``traditionally protected by our society''; tradition protected the ``unitary family'' against an ``adulterous natural father.'' Footnote 6 would consult the most specific relevant tradition, which O'Connor and Kennedy rejected as inconsistent with \textit{Griswold} and \textit{Loving}. Stevens assumed a liberty interest but found \S~4601's visitation avenue adequate. Brennan argued that biology plus a developed relationship creates a liberty interest and that Footnote 6 makes the Constitution backward-looking; White stressed that Michael satisfied \textit{Lehr}.

\section{Examples: Future Applications}
\textbf{Same-Side:} a former affair partner cannot override an intact, cohabiting marriage that wishes to raise the child. \textbf{Opposite-Side:} an unwed father whose children's mother was unmarried is protected by \textit{Stanley}. \textbf{Fence-Sitter:} a known donor invited into a married couple's family who is later cut off.

\section{Critique}
Defenders say concrete tradition keeps judges from imposing their preferences. Laurence Tribe and Michael Dorf answered that choosing a level of generality is never value-neutral, and family-law scholars fault the plurality for favoring marital form over the child's real attachments.

\section{Key Quotations}
\begin{quote}
``As we view them, they rest not upon such isolated factors but upon the historic respect --- indeed, sanctity would not be too strong a term --- traditionally accorded to the relationships that develop within the unitary family.''
\end{quote}
\noindent --- Justice Scalia, plurality opinion, 491 U.S. at 123.

\end{document}
